% Строчки для суммарной таблицы, в следующей строке прописаны используемые экземпляры узлов
%===s control, LLN0, CBCSWI1, ESCSWI1, DSCSWI1, CBCILO1, ESCILO1, DSCILO1

\begin{flushleft}
\needspace{4\baselineskip}
\color{uniblue}{\subsection{\raggedright{Контроллер присоединения (КП)}}}
\color{black}
\nopagebreak
\vspace{3mm}
\end{flushleft}

%%%%%%%%%%%%%%%%%%%%%%%%%%%%%%%%%%%%%%%%%%%%%%%%%%%%%%%%%%%%%%%%%%%%%%  ОБЩИЕ СВЕДЕНИЯ %%%%%%%%%%%%%%%%%%%%%%%%%%%%%%%%%%%%%%%%%%%%%%%%%%

\begin{flushleft}
\color{uniblue}{\subsubsection{Общие сведения}}
\color{black}
\nopagebreak
\end{flushleft}

Функциональный блок <<КП>> предназначен:
\begin{itemize}
\item для оперативного управления коммутационными аппаратами; 
\item для формирования сигнала положения коммутационного аппарата для целей АСУ ТП; 
\item для реализации оперативной блокировки управляемых коммутационных аппаратов.
\end{itemize}

В функциональном блоке <<КП>> в составе устройства заложена возможность оперативного управления коммутационными аппаратами в количестве до трех штук: 
\begin{itemize}
\item силовой выключатель (В);
\item заземляющий нож (ЗН);
\item разъединитель (Р).
\end{itemize}

\par
Ввод функционального блока <<КП>> в работу на этапе параметрирования устройства осуществляется программным переключателем SGF1 <<Ввод\_функции>>. Введенное состояние функционального блока соответствует активному состоянию сигнала <<КП: Ввод>>. Функциональный блок <<КП>> может быть выведен из работы оперативно путем активации сигнала <<ОВ КП>>, о чем сигнализирует активное состояние сигнала <<КП: Оперативный вывод>>.

\par
Логика управления режимами работы функционального блока <<КП>> выполнена в соответствии с рисунком \ref{kp:img_gen}. В таблице \ref{kp:tbl_gen} приведены параметры, необходимые для настройки функционального блока <<КП>>. Логика формирования общего сигнала переключения в составе функционального блока <<КП>> выполнена в соответствии с рисунком \ref{kp:img_per}.

\begin{longtable}{|>{\centering\arraybackslash}m{5cm}|>{\centering\arraybackslash}m{4cm}|>{\centering\arraybackslash}m{4cm}|>{\centering\arraybackslash}m{2cm}|>{\centering\arraybackslash}m{1cm}|}
\caption{Параметры для настройки функционального блока <<КП>>\hfill\vspace{-0.5\baselineskip}}\label{kp:tbl_gen}\\ 
\hline
\rowcolor{gray!30}
Параметр (Параметр на ИЧМ) & Условное обозначение на схеме & Значение/ Диапазон & Единица измерения & Шаг \\ 
\hline
\endfirsthead
\caption{\emph{Продолжение\hfill\vspace{-0.5\baselineskip}}} \\ % Отступ обозначения таблицы от таблицы и выравнивание надписи слева
\hline
\rowcolor{gray!30}
Параметр (Параметр на ИЧМ) & Условное обозначение на схеме & Значение/ Диапазон & Единица измерения & Шаг \\ 
%\hline
\endhead
%\hline
\endfoot
%\hline
\endlastfoot
%===t1*LLN0
\centering Ввод функции в работу (Ввод\_функции) & \centering SGF1 & \centering 0 = Не предусмотрено\\1 = Предусмотрено & \centering -- & \centering\arraybackslash -- \\
\hline
\centering Ввод ОБР КА в работу  (ОБР\_КА) & \centering SG1* & \centering 0 = Не предусмотрено\\1 = Предусмотрено & \centering -- & \centering\arraybackslash -- \\
\hline
%===t1
\multicolumn{5}{|l|}{* -- общий переключатель ввода оперативной блокировки КА присоединения, см. рис. \ref{kp:img_obrzn}, \ref{kp:img_obrr}} \\ 
\hline
\end{longtable}

\vspace{3mm}
\begin{figure}[H]
\centering

\includegraphics[width=0.9\textwidth,height=0.9\textheight,keepaspectratio]{\fbpath/32. КП/_latex/img/1.pdf}
\captionof{figure}{Функциональная схема управления режимами работы <<КП>>}
\label{kp:img_gen}
\end{figure}

\vspace{1mm}
\begin{figure}[H]
\centering
\includegraphics[width=0.6\textwidth,height=0.6\textheight,keepaspectratio]{\fbpath/32. КП/_latex/img/2.pdf}
\captionof{figure}{Функциональная схема формирования общего сигнала переключения коммутационного аппарата}
\label{kp:img_per}
\end{figure}

%%%%%%%%%%%%%%%%%%%%%%%%%%%%%%%%%%%%%%%%%%%%%%%%%%%%%%%%%%%%%%%%%%%%%%  УПРАВЛЕНИЕ ВЫКЛЮЧАТЕЛЕМ %%%%%%%%%%%%%%%%%%%%%%%%%%%%%%%%%%%%%%%%%%%%%%%%%%

\begin{flushleft}
\needspace{3\baselineskip}
\color{uniblue}{\subsubsection{Управление выключателем (УВ)}}
\color{black}
\nopagebreak
\end{flushleft}

\par
Ввод функции <<УВ>> в работу на этапе параметрирования устройства осуществляется программным переключателем SGF1 <<Ввод\_функции>>. Введенное состояние функции соответствует активному состоянию сигнала <<УВ: Ввод>>.
Функция <<УВ>> формирует команды на отключение и включение выключателя при оперативном управлении.

\par
Команда \emph{оперативного отключения} <<УВ: Отключить>> формируется в следующих случаях:
\begin{itemize}
	\item при появлении команды на оперативное отключение выключателя (<<Отключить В~от~ПУ>> или <<Откл. В~ИЧМ>>) и выбранным управлением по месту (сигнал <<Местное>> находится в активном состоянии);
	\item при появлении команды на оперативное отключение выключателя (<<Отключить В~от~ТУ>> или <<Откл. В~АСУ>>) и выбранным дистанционным управлением (сигнал <<Дистанционное>> находится в активном состоянии).
\end{itemize}
Прохождение команды на отключение блокируется при наличии следующих сигналов:
\begin{itemize}
	\item запрет отключения от функции контроля силового выключателя <<КСВ: Блокировка откл>>;
	\item отключенное состояние силового выключателя <<В: Отключен>>;
	\item превышение времени переключения силового выключателя <<УВ: Прев. времени~пер.>>;
	\item местное управление силовым выключателем (положение переключателя <<М/Д>> в приводе выключателя) <<В:~Ключ М/Д привода>>.
\end{itemize}

\par
Команда \emph{оперативного включения} <<УВ: Включить>> формируется в следующих случаях:
\begin{itemize}
	\item при появлении команды на оперативное включение выключателя (<<Включить В~от~ПУ>> или <<Вкл. В~ИЧМ>>) и выбранным управлением по месту (сигнал <<Местное>> находится в активном состоянии);
	\item при появлении команды на оперативное включение выключателя (<<Включить В~от~ТУ>> или <<Вкл. В~АСУ>>) и выбранным дистанционным управлением (сигнал <<Дистанционное>> находится в активном состоянии).
\end{itemize}
Прохождение команды на включение блокируется при наличии следующих сигналов:
\begin{itemize}
	\item запрет включения от функции контроля силового выключателя <<КСВ: Блокировка вкл>>;
	\item включенное состояние силового выключателя <<В:~Включен>>;
	\item превышение времени переключения силового выключателя <<УВ: Прев. времени~пер.>>;
	\item местное управление силовым выключателем (положение переключателя <<М/Д>> в приводе выключателя) <<В:~Ключ М/Д привода>>;
	\item появление сигнала аварийного отключения силового выключателя <<КСВ: В аварийно отключен>> (прохождение сигнала определяется состоянием программного переключателя SGF2 <<Бл\_вкл\_от\_авар\_откл>>).
\end{itemize}

\par
При наличии сигналов <<УВ: Включить>> и <<УВ: Отключить>> длительностью больше, чем время уставки <<T1>> формируется сигнал превышения допустимого времени переключения -- <<УВ: Прев. времени пер.>>. Активное состояние операции включения или отключения сигнализируется выходным сигналом <<УВ: Идет переключение>>.

\par
Логика функции <<УВ>> формирует сигналы состояния выключателя <<УВ: Отключено>>, <<УВ: Включено>>, <<УВ: Не определено>> и <<УВ: Неисправность>>, при этом сигналы <<УВ: Неисправность>> и <<УВ: Не определено>> выдаются с выдержкой времени, задаваемой уставкой Т2 <<Тблк>>.

\par
Алгоритм работы <<УВ>> выполнен в соответствии с рисунком \ref{kp:img_uv}. В таблице \ref{kp:tbl_uv} приведены параметры, необходимые для настройки функции <<УВ>>.

\vspace{3mm}
\begin{figure}[H]
\centering
\includegraphics[width=0.9\textwidth,height=0.9\textheight,keepaspectratio]{\fbpath/32. КП/_latex/img/uv.pdf}
\captionof{figure}{Функциональная схема <<УВ>>}
\label{kp:img_uv}
\end{figure}

\begin{longtable}{|>{\centering\arraybackslash}m{5cm}|>{\centering\arraybackslash}m{4cm}|>{\centering\arraybackslash}m{4cm}|>{\centering\arraybackslash}m{2cm}|>{\centering\arraybackslash}m{1cm}|}
\caption{Параметры для настройки функции <<УВ>>\hfill\vspace{-0.5\baselineskip}}\label{kp:tbl_uv}\\ 
\hline
\rowcolor{gray!30}
Параметр (Параметр на ИЧМ) & Условное обозначение на схеме & Значение/ Диапазон & Единица измерения & Шаг \\ 
\hline
\endfirsthead
\caption{\emph{Продолжение\hfill\vspace{-0.5\baselineskip}}} \\ % Отступ обозначения таблицы от таблицы и выравнивание надписи слева
\hline
\rowcolor{gray!30}
Параметр (Параметр на ИЧМ) & Условное обозначение на схеме & Значение/ Диапазон & Единица измерения & Шаг \\ 
%\hline
\endhead
%\hline
\endfoot
%\hline
\endlastfoot
%===t1*CBCSWI1
\centering Ввод функции в работу (Ввод\_функции) & \centering SGF1 & \centering 0 = Не предусмотрено\\1 = Предусмотрено & \centering -- & \centering\arraybackslash -- \\
\hline
\centering Блокировка включения от аварийного   отключения В (Бл\_вкл\_от\_авар\_откл) & \centering SGF2 & \centering 0 = Не предусмотрено\\1 = Предусмотрено & \centering -- & \centering\arraybackslash -- \\
\hline
\centering Допустимое время переключения выключателя (Тперекл) & \centering T1 & \centering 0,01 ... 10,00 & \centering с & \centering\arraybackslash 0,01 \\
\hline
\centering Выдержка времени подавления выдачи положения <<Не  определено>> (Тблк) & \centering T2 & \centering 0,01 ... 10,00 & \centering с & \centering\arraybackslash 0,01 \\
\hline
%===t1
\end{longtable}

%%%%%%%%%%%%%%%%%%%%%%%%%%%%%%%%%%%%%%%%%%%%%%%%%%%%%%%%%%%%%%%%%%%%%%  УПРАВЛЕНИЕ ЗН %%%%%%%%%%%%%%%%%%%%%%%%%%%%%%%%%%%%%%%%%%%%%%%%%%
\begin{flushleft}
\needspace{3\baselineskip}
\color{uniblue}{\subsubsection{Управление заземляющим ножом (УЗН)}}
\color{black}
\nopagebreak
\end{flushleft}

\par
Ввод функции <<УЗН>> в работу на этапе параметрирования устройства осуществляется программным переключателем SGF1 <<Ввод\_функции>>. Введенное состояние функции соответствует активному состоянию сигнала <<УЗН: Ввод>>.
Функция <<УЗН>> формирует команды на отключение и включение заземляющего разъединителя.

\par
Команда \emph{оперативного отключения} формируется в следующих случаях:
\begin{itemize}
	\item при появлении команды на оперативное отключение заземляющего разъединителя (<<Откл. ЗН~ИЧМ>>), выбранным управлением по месту (сигнал <<Местное>> находится в активном состоянии) и контролем разрешающего сигнала от схемы оперативной блокировки <<ОБР ЗН: Разрешить управ.>>;
	\item при появлении команды на оперативное отключение заземляющего разъединителя (<<Откл. ЗН~АСУ>>), выбранным дистанционным управлением (сигнал <<Дистанционное>> находится в активном состоянии) и контролем разрешающего сигнала от схемы оперативной блокировки <<ОБР ЗН: Разрешить управ.>>.
\end{itemize}
Прохождение команды на отключение блокируется при наличии следующих сигналов:
\begin{itemize}
	\item отключенное состояние заземляющего разъединителя <<ЗН: Отключен (б/к)>>;
	\item превышение времени переключения заземляющего разъединителя <<УЗН: Прев. времени пер.>>.
\end{itemize}

\par
Команда \emph{оперативного включения} формируется в следующих случаях:
\begin{itemize}
	\item при появлении команды на оперативное включение заземляющего разъединителя (<<Вкл. ЗН~ИЧМ>>), выбранным управлением по месту (сигнал <<Местное>> находится в активном состоянии) и контролем разрешающего сигнала от схемы оперативной блокировки <<ОБР ЗН: Разрешить управ.>>;
	\item при появлении команды на оперативное включение заземляющего разъединителя (<<Вкл. ЗН~АСУ>>), выбранным дистанционным управлением (сигнал <<Дистанционное>> находится в активном состоянии) и контролем разрешающего сигнала от схемы оперативной блокировки <<ОБР ЗН: Разрешить управ.>>.
\end{itemize}
Прохождение команды на включение блокируется при наличии следующих сигналов:
\begin{itemize}
	\item включенное состояние заземляющего разъединителя <<ЗН: Включен>>;
	\item превышение времени переключения заземляющего разъединителя <<УЗН: Прев. времени пер.>>.
\end{itemize}

\par
При наличии сигналов <<УЗН: Включить>> и <<УЗН: Отключить>> длительностью больше, чем время уставки <<T1>> формируется сигнал превышения допустимого времени переключения -- <<УЗН: Прев. времени пер.>>. Активное состояние операции включения или отключения сигнализируется выходным сигналом <<УЗН: Идет переключение>>.

\par
Логика функции <<УЗН>> формирует сигналы состояния заземляющего разъединителя <<УЗН: Отключено>>, <<УЗН: Включено>>, <<УЗН: Не определено>> и <<УЗН: Неисправность>>, при этом сигналы <<УЗН: Неисправность>> и <<УЗН: Не определено>> выдаются с выдержкой времени, задаваемой уставкой Т2 <<Тблк>>.

\par
Алгоритм работы функции <<УЗН>> выполнен в соответствии с рисунком \ref{kp:img_uzn}. В таблице \ref{kp:tbl_uzn} приведены параметры, необходимые для настройки функции <<УЗН>>.

\vspace{3mm}
\begin{figure}[H]
\centering
\includegraphics[width=0.9\textwidth,height=0.9\textheight,keepaspectratio]{\fbpath/32. КП/_latex/img/uzn.pdf}
\captionof{figure}{Функциональная схема <<УЗН>>}
\label{kp:img_uzn}
\end{figure}

\begin{longtable}{|>{\centering\arraybackslash}m{5cm}|>{\centering\arraybackslash}m{4cm}|>{\centering\arraybackslash}m{4cm}|>{\centering\arraybackslash}m{2cm}|>{\centering\arraybackslash}m{1cm}|}
\caption{Параметры для настройки функции <<УЗН>>\hfill\vspace{-0.5\baselineskip}}\label{kp:tbl_uzn}\\ 
\hline
\rowcolor{gray!30}
Параметр (Параметр на ИЧМ) & Условное обозначение на схеме & Значение/ Диапазон & Единица измерения & Шаг \\ 
\hline
\endfirsthead
\caption{\emph{Продолжение\hfill\vspace{-0.5\baselineskip}}} \\ % Отступ обозначения таблицы от таблицы и выравнивание надписи слева
\hline
\rowcolor{gray!30}
Параметр (Параметр на ИЧМ) & Условное обозначение на схеме & Значение/ Диапазон & Единица измерения & Шаг \\ 
%\hline
\endhead
%\hline
\endfoot
%\hline
\endlastfoot
%===t1*ESCSWI1
\centering Ввод функции в работу (Ввод\_функции) & \centering SGF1 & \centering 0 = Не предусмотрено\\1 = Предусмотрено & \centering -- & \centering\arraybackslash -- \\
\hline
\centering Допустимое время переключения (Тперекл) & \centering T1 & \centering 0,05 ... 100,00 & \centering с & \centering\arraybackslash 0,01 \\
\hline
\centering Выдержка времени подавления выдачи положения <<Не  определено>> (Тблк) & \centering T2 & \centering 0,05 ... 100,00 & \centering с & \centering\arraybackslash 0,01 \\
\hline
%===t1
\end{longtable}

%%%%%%%%%%%%%%%%%%%%%%%%%%%%%%%%%%%%%%%%%%%%%%%%%%%%%%%%%%%%%%%%%%%  УПРАВЛЕНИЕ РАЗЪЕДИНИТЕЛЕМ %%%%%%%%%%%%%%%%%%%%%%%%%%%%%%%%%%%%%%%%%%%%%%%%%%

\begin{flushleft}
\needspace{3\baselineskip}
\color{uniblue}{\subsubsection{Управление разъединителем (УР)}}
\color{black}
\nopagebreak
\end{flushleft}

\par
Ввод функции <<УР>> в работу на этапе параметрирования устройства осуществляется программным переключателем SGF1 <<Ввод\_функции>>. Введенное состояние функции соответствует активному состоянию сигнала <<УР: Ввод>>.
Функция <<УР>> формирует команды на отключение и включение разъединителя.

\par
Команда \emph{оперативного отключения} формируется в следующих случаях:
\begin{itemize}
	\item при появлении команды на оперативное отключениеразъединителя (<<Откл. Р~ИЧМ>>), выбранным управлением по месту (сигнал <<Местное>> находится в активном состоянии) и контролем разрешающего сигнала от схемы оперативной блокировки <<ОБР Р: Разрешить управ.>>;
	\item при появлении команды на оперативное отключение разъединителя (<<Откл. Р~АСУ>>), выбранным дистанционным управлением (сигнал <<Дистанционное>> находится в активном состоянии) и контролем разрешающего сигнала от схемы оперативной блокировки <<ОБР Р: Разрешить управ.>>.
\end{itemize}
Прохождение команды на отключение блокируется при наличии следующих сигналов:
\begin{itemize}
	\item отключенное состояние разъединителя <<Р: Отключен (б/к)>>;
	\item превышение времени переключения разъединителя <<УР: Прев. времени пер.>>.
\end{itemize}

\par
Команда \emph{оперативного включения} формируется в следующих случаях:
\begin{itemize}
	\item при появлении команды на оперативное включение разъединителя (<<Вкл. Р~ИЧМ>>), выбранным управлением по месту (сигнал <<Местное>> находится в активном состоянии) и контролем разрешающего сигнала от схемы оперативной блокировки <<ОБР Р: Разрешить управ.>>;
	\item при появлении команды на оперативное включение разъединителя (<<Вкл. Р~АСУ>>), выбранным дистанционным управлением (сигнал <<Дистанционное>> находится в активном состоянии) и контролем разрешающего сигнала от схемы оперативной блокировки <<ОБР Р: Разрешить управ.>>.
\end{itemize}
Прохождение команды на включение блокируется при наличии следующих сигналов:
\begin{itemize}
	\item включенное состояние разъединителя <<Р: Включен>>;
	\item превышение времени переключения разъединителя <<УР: Прев. времени пер.>>.
\end{itemize}

\par
При наличии сигналов <<УР: Включить>> и <<УР: Отключить>> длительностью больше, чем время уставки <<T1>> формируется сигнал превышения допустимого времени переключения -- <<УР: Прев. времени пер.>>. Активное состояние операции включения или отключения сигнализируется выходным сигналом <<УР: Идет переключение>>.

\par
Логика функции <<УР>> формирует сигналы состояния разъединителя <<УР: Отключено>>, <<УР: Включено>>, <<УР: Не определено>> и <<УР: Неисправность>>, при этом сигналы <<УР: Неисправность>> и <<УР: Не определено>> выдаются с выдержкой времени, задаваемой уставкой Т2 <<Тблк>>.

\par
Алгоритм работы функции <<УР>> выполнен в соответствии с рисунком \ref{kp:img_ur}. В таблице \ref{kp:tbl_ur} приведены параметры, необходимые для настройки функции <<УР>>. 

\vspace{3mm}
\begin{figure}[H]
\centering
\includegraphics[width=0.9\textwidth,height=0.9\textheight,keepaspectratio]{\fbpath/32. КП/_latex/img/ur.pdf}
\captionof{figure}{Функциональная схема <<УР>>}
\label{kp:img_ur}
\end{figure}

\begin{longtable}{|>{\centering\arraybackslash}m{5cm}|>{\centering\arraybackslash}m{4cm}|>{\centering\arraybackslash}m{4cm}|>{\centering\arraybackslash}m{2cm}|>{\centering\arraybackslash}m{1cm}|}
\caption{Параметры для настройки функции <<УР>>\hfill\vspace{-0.5\baselineskip}}\label{kp:tbl_ur}\\ 
\hline
\rowcolor{gray!30}
Параметр (Параметр на ИЧМ) & Условное обозначение на схеме & Значение/ Диапазон & Единица измерения & Шаг \\ 
\hline
\endfirsthead
\caption{\emph{Продолжение\hfill\vspace{-0.5\baselineskip}}} \\ % Отступ обозначения таблицы от таблицы и выравнивание надписи слева
\hline
\rowcolor{gray!30}
Параметр (Параметр на ИЧМ) & Условное обозначение на схеме & Значение/ Диапазон & Единица измерения & Шаг \\ 
%\hline
\endhead
%\hline
\endfoot
%\hline
\endlastfoot
%===t1*DSCSWI1
\centering Ввод функции в работу (Ввод\_функции) & \centering SGF1 & \centering 0 = Не предусмотрено\\1 = Предусмотрено & \centering -- & \centering\arraybackslash -- \\
\hline
\centering Допустимое время переключения (Тперекл) & \centering T1 & \centering 0,05 ... 100,00 & \centering с & \centering\arraybackslash 0,01 \\
\hline
\centering Выдержка времени подавления выдачи положения <<Не  определено>> (Тблк) & \centering T2 & \centering 0,05 ... 100,00 & \centering с & \centering\arraybackslash 0,01 \\
\hline
%===t1
\end{longtable}

%%%%%%%%%%%%%%%%%%%%%%%%%%%%%%%%%%%%%%%%%%%%%%%%%%%%%%%%%%%%%%%%%%%%%%  ОПЕРАТИВНАЯ БЛОКИРОВКА В %%%%%%%%%%%%%%%%%%%%%%%%%%%%%%%%%%%%%%%%%%%%%%%%%%

\begin{flushleft}
\needspace{3\baselineskip}
\color{uniblue}{\subsubsection{Оперативная блокировка выключателя (ОБР В)}}
\color{black}
\nopagebreak
\end{flushleft}

\par
Функция оперативной и технологической блокировки силового выключателя вводится в работу при введенной функции <<УВ>>.
Выполнение операций с силовым выключателем возможно при отсутствии сигнала о текущем выполнении каких-либо переключений на присоединении (определяется программным переключателем SGF1 <<ОБР\_В>>).

\par
Алгоритм работы <<ОБР В>> выполнен в соответствии с рисунком \ref{kp:img_obrv}. В таблице \ref{kp:tbl_obrv} приведены параметры, необходимые для настройки функции <<ОБР В>>.

\vspace{3mm}
\begin{figure}[H]
\centering
\includegraphics[width=0.7\textwidth,height=0.7\textheight,keepaspectratio]{\fbpath/32. КП/_latex/img/obrv.pdf}
\captionof{figure}{Функциональная схема <<ОБР В>>}
\label{kp:img_obrv}
\end{figure}

\begin{longtable}{|>{\centering\arraybackslash}m{5cm}|>{\centering\arraybackslash}m{4cm}|>{\centering\arraybackslash}m{4cm}|>{\centering\arraybackslash}m{2cm}|>{\centering\arraybackslash}m{1cm}|}
\caption{Параметры для настройки <<ОБР В>>\hfill\vspace{-0.5\baselineskip}}\label{kp:tbl_obrv}\\ 
\hline
\rowcolor{gray!30}
Параметр (Параметр на ИЧМ) & Условное обозначение на схеме & Значение/ Диапазон & Единица измерения & Шаг \\ 
\hline
\endfirsthead
\caption{\emph{Продолжение\hfill\vspace{-0.5\baselineskip}}} \\ % Отступ обозначения таблицы от таблицы и выравнивание надписи слева
\hline
\rowcolor{gray!30}
Параметр (Параметр на ИЧМ) & Условное обозначение на схеме & Значение/ Диапазон & Единица измерения & Шаг \\ 
%\hline
\endhead
%\hline
\endfoot
%\hline
\endlastfoot
%===t1*CBCILO1
\centering Ввод функции в работу (Ввод\_функции) & \centering SGF1 & \centering 0 = Не предусмотрено\\1 = Предусмотрено & \centering -- & \centering\arraybackslash -- \\
\hline
%===t1
\end{longtable}

%%%%%%%%%%%%%%%%%%%%%%%%%%%%%%%%%%%%%%%%%%%%%%%%%%%%%%%%%%%%%%%%%%%%%%  ОПЕРАТИВНАЯ БЛОКИРОВКА ЗН %%%%%%%%%%%%%%%%%%%%%%%%%%%%%%%%%%%%%%%%%%%%%%%%%%

\begin{flushleft}
\needspace{3\baselineskip}
\color{uniblue}{\subsubsection{Оперативная блокировка ЗН (ОБР ЗН)}}
\color{black}
\nopagebreak
\end{flushleft}

\par
Ввод <<ОБР ЗН>> в работу на этапе параметрирования устройства осуществляется общим программным переключателем SG1 <<ОБР\_КА>> (см. таблицу \ref{kp:tbl_gen}).

\par
Принудительное деблокирование управления заземляющего разъединителя трансформатора стороны ВН от оперативной блокировки осуществляется активацией входного сигнала <<Программ. деблок. ЗН>>.
\par
Разрешение на отключение заземляющего разъединителя выдается без дополнительных условий.
\par
Разрешение на включение заземляющего разъединителя выдается при выполнении необходимых условий:
\begin{itemize}
	\item разъединитель трансформатора стороны ВН отключен; 
	\item разъединитель трансформатора стороны НН1 отключен либо выкатной элемент выключателя НН1 трансформатора в выкаченном состоянии (контроль сигнала определяется состоянием программного переключателя SGF1 <<Контр\_КА3>>); 
	\item разъединитель трансформатора стороны НН2 отключен либо выкатной элемент выключателя НН2 трансформатора в выкаченном состоянии (контроль сигнала определяется состоянием программного переключателя SGF2 <<Контр\_КА4>>).
\end{itemize}

\par
Алгоритм работы <<ОБР ЗН>> выполнен в соответствии с рисунком \ref{kp:img_obrzn}. В таблице \ref{kp:tbl_obrzn} приведены параметры, необходимые для настройки функции <<ОБР ЗН>>.

\vspace{3mm}
\begin{figure}[H]
\centering

\includegraphics[width=0.8\textwidth,height=0.8\textheight,keepaspectratio]{\fbpath/32. КП/_latex/img/obrzn.pdf}
\captionof{figure}{Функциональная схема <<ОБР ЗН>>}
\label{kp:img_obrzn}
\end{figure}

\begin{longtable}{|>{\centering\arraybackslash}m{5cm}|>{\centering\arraybackslash}m{4cm}|>{\centering\arraybackslash}m{4cm}|>{\centering\arraybackslash}m{2cm}|>{\centering\arraybackslash}m{1cm}|}
\caption{Параметры для настройки функции <<ОБР ЗН>>\hfill\vspace{-0.5\baselineskip}}\label{kp:tbl_obrzn}\\ 
\hline
\rowcolor{gray!30}
Параметр (Параметр на ИЧМ) & Условное обозначение на схеме & Значение/ Диапазон & Единица измерения & Шаг \\ 
\hline
\endfirsthead
\caption{\emph{Продолжение\hfill\vspace{-0.5\baselineskip}}} \\ % Отступ обозначения таблицы от таблицы и выравнивание надписи слева
\hline
\rowcolor{gray!30}
Параметр (Параметр на ИЧМ) & Условное обозначение на схеме & Значение/ Диапазон & Единица измерения & Шаг \\ 
%\hline
\endhead
%\hline
\endfoot
%\hline
\endlastfoot
%===t1*ESCILO1
\centering  Контроль КА3 (Контр\_КА3) & \centering SGF1 & \centering 0 = Не предусмотрено\\1 = Предусмотрено & \centering -- & \centering\arraybackslash -- \\
\hline
\centering  Контроль КА4 (Контр\_КА4) & \centering SGF2 & \centering 0 = Не предусмотрено\\1 = Предусмотрено & \centering -- & \centering\arraybackslash -- \\
\hline
%===t1
\end{longtable}

%%%%%%%%%%%%%%%%%%%%%%%%%%%%%%%%%%%%%%%%%%%%%%%%%%%%%%%%%%%%%%%% ОПЕРАТИВНАЯ БЛОКИРОВКА Р %%%%%%%%%%%%%%%%%%%%%%%%%%%%%%%%%%%%%%%%%%%%%%%%%%%%%%%%%%%%%%%%%

\begin{flushleft}
\needspace{2\baselineskip}
\color{uniblue}{\subsubsection{Оперативная блокировка разъединителя (ОБР Р)}}
\color{black}
\nopagebreak
\end{flushleft}

\par
Ввод <<ОБР Р>> в работу на этапе параметрирования устройства осуществляется общим программным переключателем SG1 <<ОБР\_КА>> (см. таблицу \ref{kp:tbl_gen}).

\par
Принудительное деблокирование управления разъединителя трансформатора стороны ВН от оперативной блокировки осуществляется активацией входного сигнала <<Программ. деблок. Р>>.
\par
Разрешение на отключение разъединителя выдается без дополнительных условий.
\par
Разрешение на включение разъединителя выдается при выполнении необходимых условий:
\begin{itemize}
	\item заземляющий разъединитель трансформатора стороны ВН отключен; 
	\item заземляющий разъединитель трансформатора стороны НН1 отключен (контроль сигнала определяется состоянием программного переключателя SGF1 <<Контр\_КА1>>); 
	\item заземляющий разъединитель трансформатора стороны НН2 отключен (контроль сигнала определяется состоянием программного переключателя SGF2 <<Контр\_КА2>>).
\end{itemize}

\par
Алгоритм работы <<ОБР Р>> выполнен в соответствии с рисунком \ref{kp:img_obrr}. В таблице \ref{kp:tbl_obrr} приведены параметры, необходимые для настройки функции <<ОБР Р>>.

\vspace{3mm}
\begin{figure}[H]
\centering
\includegraphics[width=0.8\textwidth,height=0.8\textheight,keepaspectratio]{\fbpath/32. КП/_latex/img/obrr.pdf}
\captionof{figure}{Функциональная схема <<ОБР Р>>}
\label{kp:img_obrr}
\end{figure}

\begin{longtable}{|>{\centering\arraybackslash}m{5cm}|>{\centering\arraybackslash}m{4cm}|>{\centering\arraybackslash}m{4cm}|>{\centering\arraybackslash}m{2cm}|>{\centering\arraybackslash}m{1cm}|}
\caption{Параметры для настройки функции <<ОБР Р>>\hfill\vspace{-0.5\baselineskip}}\label{kp:tbl_obrr}\\ 
\hline
\rowcolor{gray!30}
Параметр (Параметр на ИЧМ) & Условное обозначение на схеме & Значение/ Диапазон & Единица измерения & Шаг \\ 
\hline
\endfirsthead
\caption{\emph{Продолжение\hfill\vspace{-0.5\baselineskip}}} \\ % Отступ обозначения таблицы от таблицы и выравнивание надписи слева
\hline
\rowcolor{gray!30}
Параметр (Параметр на ИЧМ) & Условное обозначение на схеме & Значение/ Диапазон & Единица измерения & Шаг \\ 
%\hline
\endhead
%\hline
\endfoot
%\hline
\endlastfoot
%===t1*DSCILO1
\centering  Контроль КА1 (Контр\_КА1) & \centering SGF1 & \centering 0 = Не предусмотрено\\1 = Предусмотрено & \centering -- & \centering\arraybackslash -- \\
\hline
\centering  Контроль КА2 (Контр\_КА2) & \centering SGF2 & \centering 0 = Не предусмотрено\\1 = Предусмотрено & \centering -- & \centering\arraybackslash -- \\
\hline
%===t1
\end{longtable}
